\section{Conclusions}
\label{sec:conclusions}

We presented a three-stage funnel that extracts author-stated research gaps as unmodified source sentences. Rules remove unlikely text, a zero-shot NLI classifier selects candidates without gap-labelled training data, and an LLM filter removes false positives in one of three modes set by a single instruction.

The stages have distinct roles. Stage~A helps when sections are informative, Stage~B sets the recall ceiling, and Stage~C supplies most of the precision gain. The modes shift precision and recall in the intended direction on every dataset. In high-precision mode, precision is 66.2\% at 84.9\% recall on SAL-40 and 76.5\% at 75.5\% recall on FWS. On our gold corpus, whose labels come from LLMs and are not yet checked by people, 86.6\% of kept sentences match those labels at 57.4\% recall. Some gold test candidates influenced funnel design, so these results are provisional.

With the 27B LLM, the funnel has higher observed $F_1$ than whole-paper prompting on SAL-40, but the confidence interval includes zero and only the funnel receives section labels. It also has higher $F_1$ on FWS, but a broader baseline prompt favours the funnel. Whole-paper prompting has higher observed $F_1$ on arXiv papers against labels made from whole-paper inputs. With the 4B LLM, the funnel is substantially better on long arXiv papers. Within the two tested Qwen models, filtering helps most for the 4B model on long inputs; model size is confounded with model generation.

For Stage~B, the training-free NLI classifier comes within 0.021 $F_1$ of a supervised PubMedBERT in high-precision mode on SAL, but is 8.4 points less precise. Fine-tuning on LLM labels improves agreement with those labels, not with the human-labelled benchmarks.

The next steps are the human validation of the gold corpus, comparisons with published extraction methods, and the use of the extracted gaps for automated research-idea generation.
